\section{The exact Fock normalizer and equality at large primes}
\label{fh-section}

We now compute the normalizing integer $C_\nu$ of
\eqref{coef:fock}. This gives an explicit integer which kills
$\nabla_\nu(\mathbb Z)/\Delta_\nu(\mathbb Z)$, and proves the
local equality of these lattices needed at the level primes.
The source obtains that equality from simplicity in the lowest
alcove.\cite{alpo-uniform-source} The calculation below instead
uses polynomial differential operators and the integral duality
already proved in Theorem~\ref{coef:integral-duality}.

\subsection{A differential determinant identity}

Let $X=(x_{ia})_{1\le i,a\le h}$ be a square matrix of variables,
write $\partial_{ia}=\partial/\partial x_{ia}$, and put
$E_{ij}=\sum_a x_{ia}\partial_{ja}$.
Composition of operators acts from right to left; in all products
below the factors are written from left to right in increasing $j$.

\begin{lemma}[The differential Capelli identity]\label{fh-capelli}
On $\mathbb Z[X]$ one has
\begin{equation}\label{fh-capelli-identity}
 \det(X)\det(\partial)=
 \sum_{\sigma\in\mathfrak S_h}\operatorname{sgn}(\sigma)
       \prod_{j=1}^h
           \bigl(E_{\sigma(j),j}+(h-j)\delta_{\sigma(j),j}\bigr).
\end{equation}
On the left, multiplication by $\det(X)$ follows differentiation.
\end{lemma}
\begin{proof}
Adjoin exterior generators $e_1,\ldots,e_h$, so
$e_ie_j=-e_je_i$ and $e_i^2=0$. Derivatives act only on the
polynomial variables. Set
\[
       \xi_a=\sum_i e_i x_{ia},\qquad
       A_j=\sum_a\xi_a\partial_{ja}=\sum_i e_iE_{ij}.
\]
We prove by descending induction on $j$ the operator identity
\begin{equation}\label{fh-exterior-ordering}
 \prod_{i=j}^h\bigl(A_i+(h-i)e_i\bigr)
   =\sum_{a_j,\ldots,a_h=1}^h
        \xi_{a_j}\cdots\xi_{a_h}
        \partial_{j,a_j}\cdots\partial_{h,a_h}.
\end{equation}
The empty tail is the identity. Suppose the assertion is known
for $j+1$, and denote its right side by $N_{j+1}$.
When $A_j$ acts on $N_{j+1}$, the terms in which its derivative
passes all $h-j$ exterior factors give exactly the right side
of \eqref{fh-exterior-ordering}. There is one additional term
for each factor hit by the derivative, since
\[
                         \partial_{ja}\xi_b=e_j\delta_{ab}.
\]
Each such term is $-e_jN_{j+1}$. To check the sign, if the
derivative hits the $t$th exterior factor of the tail, move the
resulting $e_j$ to the front, across $t$ factors, then move the
initial $\xi_a$ into its former position, across $t-1$ factors.
The combined sign is $(-1)^t(-1)^{t-1}=-1$.
Thus the total additional contribution is
$-(h-j)e_jN_{j+1}$, which the displayed shift cancels.
This proves the induction, using only integral operator identities.

Take the coefficient of $e_1\cdots e_h$ in
\eqref{fh-exterior-ordering} for $j=1$.
On the left it is the ordered sum on the right side of
\eqref{fh-capelli-identity}. On the right, only permutations of
$a_1,\ldots,a_h$ contribute: a repeated index makes the exterior
product zero. For $a_j=\tau(j)$ that coefficient is
$\operatorname{sgn}(\tau)\det(X)$.
The remaining sum of commuting partial derivatives is
$\det(\partial)$. This proves \eqref{fh-capelli-identity}.
\end{proof}

\subsection{The hook product}

Let $\nu$ be a partition. Its diagram consists of the pairs
$(i,j)$ with $1\le j\le\nu_i$. Write
$\nu'_j=\#\{i:\nu_i\ge j\}$ and define the hook length of
the pair $(i,j)$ by
\[
                  h_\nu(i,j)=\nu_i-j+\nu'_j-i+1.
\]
It is a positive integer: it counts that pair, the pairs to its
right in the same row, and those below it in the same column.
Zero rows and extra unused variables do not change the highest
polynomial $H_\nu$ or its Fock norm. Empty products are one.

\begin{theorem}[Exact Fock normalization]\label{fh-hook-product}
For the highest polynomial and the unnormalized Fock form of
Section~\ref{coef:section},
\begin{equation}\label{fh-hook-equality}
                   C_\nu=[H_\nu,H_\nu]
                          =\prod_{(i,j)\in\nu}h_\nu(i,j).
\end{equation}
If $\nu$ has $h>0$ positive parts, all factors on the right
are at most $\nu_1+h-1$.
\end{theorem}
\begin{proof}
For the zero partition, $H_\nu=1$ and $C_\nu=1$.
Otherwise remove the zero parts and work on the first $h$ rows
and columns. All $\nu_i$ are then positive. Define
$\nu^-=(\nu_1-1,\ldots,\nu_h-1)$, omitting any zero parts
when computing its highest polynomial. The defining products of
initial principal minors give
\[
                             H_\nu=\det(X)H_{\nu^-}.
\]
The operators $E_{ij}$ with $i<j$ kill $H_\nu$: they
differentiate addition of an earlier row to a later one, which
preserves every initial principal minor involving the later row;
the other factors do not involve that row. Also
$E_{ii}H_\nu=\nu_iH_\nu$ by row homogeneity.

In a term of the ordered determinant
\eqref{fh-capelli-identity} for a nonidentity permutation
$\sigma$, let $j$ be its greatest nonfixed index. Then
$\sigma(j)<j$. All factors to its right are diagonal and
act on $H_\nu$ by scalars, while the factor at $j$ is a
raising operator and kills it. Only the identity permutation
survives. Hence
\[
 \det(X)\det(\partial)H_\nu
                  =\left(\prod_{i=1}^h(\nu_i+h-i)\right)H_\nu.
\]
Cancellation of the nonzero polynomial $\det(X)$ yields
\begin{equation}\label{fh-lowering}
 \det(\partial)H_\nu
                  =\left(\prod_{i=1}^h(\nu_i+h-i)\right)H_{\nu^-}.
\end{equation}

For the Fock pairing, multiplication by $x_{ia}$ is adjoint to
$\partial_{ia}$; this follows immediately from the factorial
weight of a monomial. By composition and linearity,
multiplication by $\det(X)$ is adjoint to $\det(\partial)$.
Consequently \eqref{fh-lowering} gives the exact recurrence
\begin{equation}\label{fh-column-recurrence}
                 C_\nu=
                    \left(\prod_{i=1}^h(\nu_i+h-i)\right)C_{\nu^-}.
\end{equation}
The factors in this product are exactly the hook lengths in
the first column of $\nu$. Removing that column gives
$\nu^-$, and does not change the hook lengths of the remaining
pairs: their column indices decrease by one, their row lengths
decrease by one, and their column heights are the former heights
of the next columns. Induction on the number of pairs proves
\eqref{fh-hook-equality}. Finally
$h_\nu(i,j)\le\nu_1+h-1$, directly from its definition.
\end{proof}

\subsection{An explicit annihilator and the exact prime condition}

\begin{corollary}\label{fh-annihilator}
For every partition $\nu$,
\begin{equation}\label{fh-inclusions}
              C_\nu\nabla_\nu(\mathbb Z)
                    \subseteq\Delta_\nu(\mathbb Z)
                    \subseteq\nabla_\nu(\mathbb Z).
\end{equation}
In particular, if a prime $p$ divides no hook length of $\nu$,
then
\[
                \Delta_\nu(\mathbb Z_{(p)})
                          =\nabla_\nu(\mathbb Z_{(p)}).
\]
The equality remains true after extension to any
characteristic-zero coefficient domain in which $C_\nu$ is a unit.
\end{corollary}
\begin{proof}
For $F,G\in\nabla_\nu(\mathbb Z)$,
\[
                   B_\nu(C_\nu F,G)=[F,G]\in\mathbb Z.
\]
By Theorem~\ref{coef:integral-duality}, this says exactly that
$C_\nu F\in\Delta_\nu(\mathbb Z)$. The other inclusion was
part of the construction. Tensoring these inclusions with the
coefficient ring and using the perfect bases gives the assertion
whenever $C_\nu$ becomes a unit. Equation~\eqref{fh-hook-equality}
supplies the stated sufficient prime condition.
\end{proof}

Let $\lambda_1\ge\cdots\ge\lambda_N$ be a dominant integral
weight with $N\ge1$, and put $\nu_i=\lambda_i-\lambda_N$.
Then $\nu_N=0$. If $\nu$ is nonempty, its number $h$ of
positive parts is at most $N-1$, so every hook length is at most
\[
                          \lambda_1-\lambda_N+N-2.
\]
We therefore obtain the following precise conclusion.

\begin{corollary}\label{fh-large-prime-equality}
If $p$ is prime and
\begin{equation}\label{fh-prime-range}
                       \lambda_1-\lambda_N+N-1\le p,
\end{equation}
the standard Weyl and dual-Weyl lattices of weight $\lambda$
are equal after localization at $p$. The same is true at every
prime over $p$ in a number field, and after transport by a common
local basis. More explicitly, the single positive integer
$c=C_\nu$, independent of the polynomial, satisfies
\[
       p\nmid c,\qquad
       c\nabla_\nu(\mathbb Z)\subseteq\Delta_\nu(\mathbb Z).
\]
For a constant weight, including $N=1$, one can take $c=1$.
The module in rank zero also consists of constants and has $c=1$.
\end{corollary}
\begin{proof}
For a nonempty $\nu$, \eqref{fh-prime-range} makes every
positive hook length at most $p-1$. Apply
Corollary~\ref{fh-annihilator}. The determinant twist changes
the group action but not these standard polynomial lattices or
their normalized pairing. Localization, coefficient extension
and a common invertible transport preserve equality.
The constant cases follow from $H_\nu=1$.
\end{proof}

The actual source uses the stronger hypothesis
$\lambda_1-\lambda_N+N<p$: its quantity called the spread
includes $N$, not $N-1$. Thus Corollary~\ref{fh-large-prime-equality}
proves the precise local lattice assertion used there, with a
slightly larger allowed range and an explicit clearing integer.
It does not require a modular simplicity theorem or a passage
from equality of dimensions to equality of integral lattices.

The prime condition cannot be omitted in general. For
$\lambda=(2,0)$, the dual-Weyl lattice is
$\mathbb Zx^2\oplus\mathbb Zxy\oplus\mathbb Zy^2$, whereas
the Weyl lattice is
$\mathbb Zx^2\oplus\mathbb Z(2xy)\oplus\mathbb Zy^2$.
Here $C_{(2)}=2$, and the two lattices remain different at $p=2$.
This is the example of Section~\ref{coef:section}; it lies
outside \eqref{fh-prime-range}. No equality of the two arbitrary
integral lattices is being asserted.


% Integral initial-minor spanning, proved by an elementary leading-term argument.
% Labels in this fragment use the prefix sm-.
